% ============================================================
%  Chapter 5 – System Components
% ============================================================

\section{System Components}
\label{sec:system-components}

The four principal components examined in this chapter---dual-channel retrieval, query-dependent
authority scoring, the multi-stage generation pipeline, and the citation integrity system---are
each responsible for a distinct class of computational and inferential tasks. Their interaction
is not merely sequential: authority scores inform retrieval ranking, retrieval outputs constrain
generation scope, and the citation integrity system operates as a transversal invariant that
binds every generative operation to its evidential basis. The chapter analyses each component in
terms of its formal design objectives, its algorithmic mechanisms, and the trade-offs that shaped
its architecture.


%%%%%%%%%%%%%%%%%%%%%%%%%%%%%%%%%%%%%%%%%%%%%%%%%%%%
\subsection{Dual-Channel Retrieval}
\label{sec:dual-channel-retrieval}

\subsubsection{Problem Definition}

Let $\mathcal{C} = \{c_1, c_2, \ldots, c_N\}$ denote the corpus of evidence fragments, where
each chunk $c_i$ is a contiguous textual passage extracted from a parliamentary speech and
characterised by a dense vector embedding $\mathbf{e}_i \in \mathbb{R}^d$, a speaker identity
$\sigma(c_i)$, and a parliamentary group affiliation $\gamma(c_i)$. Given a natural language
query $q$ with embedding $\mathbf{q} = \phi(q) \in \mathbb{R}^d$, the retrieval problem
is to select a subset $\mathcal{R} \subseteq \mathcal{C}$ of fixed maximum cardinality that
satisfies three partially competing requirements.

The first requirement is \textit{semantic relevance}: the selected fragments should carry
informational content that is pertinent to the query's topic. The second is \textit{political
coverage}: the selected set should include substantive evidence from each of the principal
parliamentary groups $G = \{g_1, \ldots, g_M\}$, reflecting the multi-stakeholder nature of
legislative debate. Formally, for a coverage threshold $\kappa$, the constraint reads
$|\{c_i \in \mathcal{R} : \gamma(c_i) = g_j\}| \geq \kappa$ for each group $g_j \in G$ that
has produced topically relevant material. The third requirement is \textit{authority
concentration}: among fragments from the same group, the set should prefer those attributed to
speakers with demonstrated expertise on the query's domain.

These three requirements cannot be simultaneously optimised by any single retrieval mechanism.
Semantic relevance maximisation without coverage constraints tends to concentrate retrieved
fragments among the most voluminous speakers, reproducing in the retrieval stage the very
asymmetries that the system is designed to counteract. Coverage enforcement without relevance
weighting degrades the quality of the evidentiary base. Authority concentration applied
uniformly across groups introduces a different form of bias: it can suppress minority voices
within a group even when their passages are more directly pertinent to the query.

The dual-channel architecture addresses this tension by separating the retrieval function into
two complementary mechanisms that operate on different information sources---semantic content
and structural metadata---and whose outputs are subsequently merged through a ranking function
that integrates all three requirements.

\subsubsection{Dense Channel: Vector Similarity Search}

The dense retrieval channel operates on the assumption that semantic proximity in the shared
embedding space is a reliable proxy for topical relevance. Given the query embedding
$\mathbf{q}$, the channel assigns to each chunk $c_i$ a relevance score defined as the cosine
similarity between the two embeddings:

\begin{equation}
    R(c_i, q) = \frac{\mathbf{e}_i \cdot \mathbf{q}}{\|\mathbf{e}_i\| \cdot \|\mathbf{q}\|}.
\end{equation}

The channel retrieves the $K_\text{dense}$ chunks with the highest values of $R(c_i, q)$,
subject to a minimum similarity threshold $\tau_\text{dense}$ below which a fragment is
considered topically unrelated regardless of its rank. This threshold-gating prevents the set
from being populated by weakly relevant material when the corpus contains few genuinely pertinent
passages.

The cosine metric has several properties that make it appropriate for this domain. It is
invariant to the magnitude of the embedding vectors, which means that the length of a passage---and
consequently the total information it encodes---does not inflate its retrieval score. It is
also symmetric in the embedding space, so that the same geometric intuition applies whether the
query is short and underspecified or long and detailed, provided both are mapped to the same
embedding space by the shared encoder.

The primary limitation of this channel is its sensitivity to the distributional properties of
the training corpus of the embedding model. In parliamentary discourse, certain topics are
debated far more frequently by the majority coalition than by opposition groups, simply as a
consequence of agenda-setting power. A purely density-based retrieval mechanism will therefore
tend to over-represent majority speakers on most policy topics: their speeches are more numerous,
and the cumulative effect of multiple high-similarity fragments from the same group displaces
minority voices from the top-$K$ set before any diversity constraint is applied. A second
limitation is that the dense channel has no access to the relational structure of the parliamentary
record: it cannot distinguish between a deputy who has spoken tangentially about a topic and one
who has introduced and actively sponsored legislation in the same domain.

\subsubsection{Graph Channel: Structural Traversal}

The graph retrieval channel addresses both limitations of the dense channel by exploiting the
relational metadata of the parliamentary knowledge graph. Its operating principle is that
\textit{demonstrated legislative commitment} to a policy domain is a stronger signal of
authoritativeness than incidental discursive proximity, and that this commitment is encoded in
the structural relationships between deputies, parliamentary acts, and committee memberships---none
of which is directly accessible through embedding similarity.

The channel proceeds in two traversal steps. In the first step, it identifies the set of
parliamentary acts $\mathcal{A}_q \subseteq \mathcal{A}$ whose pre-stored thematic embeddings
$\{\mathbf{a}_k\}$ are sufficiently similar to the query embedding:

\begin{equation}
    \mathcal{A}_q = \left\{ a_k \in \mathcal{A} : \frac{\mathbf{a}_k \cdot \mathbf{q}}{\|\mathbf{a}_k\| \cdot \|\mathbf{q}\|} \geq \tau_\text{graph} \right\}.
\end{equation}

In the second step, the graph is traversed from each act in $\mathcal{A}_q$ along signatorship
relationships to the set of deputies who co-sponsored it, and from those deputies to the speeches
they delivered in parliamentary sessions. The resulting evidence fragments are retrieved as the
graph channel's output.

This traversal embeds a structural prior that is absent from the dense channel: it rewards
deputies whose legislative activity---as opposed merely to their verbal contributions---places
them in proximity to the query's domain. A deputy who has sponsored multiple acts on environmental
regulation is structurally positioned as an authority on such matters regardless of whether their
debate contributions happen to be topically dense; conversely, a deputy who has spoken frequently
about the environment without ever sponsoring relevant legislation receives no structural bonus.

The complementarity of the two channels is not incidental. The dense channel is effective at
capturing topically relevant contributions that arise from debate but are not associated with
formal legislative activity---for instance, opposition critiques, expert testimonies solicited
during committee hearings, or speeches that address the consequences of legislation sponsored
by other parties. The graph channel introduces evidence from speakers who are structurally
committed to the domain but whose verbal contributions may rank below the density threshold due
to lexical variation or the relative frequency of their topic in the overall corpus. The union
of the two outputs is therefore substantively more diverse in terms of both ideological
representation and evidential type than either channel alone.

\subsubsection{Channel Merging and Ranking}
\label{sec:channel-merging}

The combined candidate pool $\mathcal{R}_0 = \mathcal{R}_\text{dense} \cup \mathcal{R}_\text{graph}$
is scored by a multi-factor ranking function that transforms the raw retrieval outputs into an
ordered evidence set suitable for generation. The ranking function integrates five dimensions of
evaluation into a composite score:

\begin{equation}
    S(c_i) = w_R \cdot R(c_i, q) + w_A \cdot A(\sigma(c_i), q) + w_D \cdot D(c_i, \mathcal{R}_0)
             + w_\Sigma \cdot \Sigma(c_i) + w_B \cdot B(c_i, \mathcal{R}_0),
\end{equation}

where $R(c_i, q)$ is the semantic relevance score defined above, $A(\sigma(c_i), q)$ is the
authority score of the speaker attributed to $c_i$ (computed by the authority scoring module
described in the following section), $D(c_i, \mathcal{R}_0)$ is a diversity penalty,
$\Sigma(c_i)$ is a salience estimate, and $B(c_i, \mathcal{R}_0)$ is a coverage balance term.
The weights $\{w_R, w_A, w_D, w_\Sigma, w_B\}$ are fixed hyperparameters of the ranking
function.

The \textit{diversity penalty} $D(c_i, \mathcal{R}_0)$ is a decreasing function of the number
of fragments already selected from the same parliamentary group as $c_i$. Its effect is to
suppress the score of fragments from over-represented groups without eliminating them from
consideration entirely: a high-authority, highly relevant fragment from a dominant group is not
excluded, but it must clear a higher effective threshold than an equivalent fragment from a
group with fewer representatives in the current selection.

The \textit{salience estimate} $\Sigma(c_i)$ reflects the attitudinal specificity and
informational density of the passage. Chunks consisting predominantly of procedural formulae,
formal acknowledgements, or parliamentary preambles carry low salience, as they convey no
substantive policy position. When such a chunk is retrieved---typically because it shares
lexical proximity with query terms that also appear in procedural contexts---a
\textit{neighbourhood expansion} mechanism replaces or augments it with the immediately succeeding
chunk in the speech, exploiting the directional continuity relation between consecutive
fragments. This expansion is predicated on the empirical observation that procedurally weak
openings in parliamentary speeches are almost invariably followed by their substantive content.

The \textit{coverage balance term} $B(c_i, \mathcal{R}_0)$ provides a corrective bonus to
fragments from groups that are under-represented in the current candidate pool relative to the
coverage threshold $\kappa$. If the final ranked selection still fails to meet the coverage
constraint after exhausting the primary candidate pool, targeted supplementary queries are
executed, restricted by group affiliation, to fill the remaining coverage gaps. These
supplementary queries inherit the same similarity threshold but operate on a narrower scope,
ensuring that the resulting additions are at least minimally relevant to the query even if they
cannot be ranked competitively against the primary candidates.

The trade-off embedded in this ranking design is between optimality in expectation and robustness
to distributional imbalances in the corpus. A ranking function that optimises only the first
three terms---relevance, authority, and diversity---will produce excellent results when the corpus
is uniformly distributed across parliamentary groups but will degrade systematically when some
groups are structurally underrepresented. The salience and coverage terms are not improvements
in the Pareto sense; they trade marginal retrieval precision for guaranteed political coverage,
a trade-off that is explicitly demanded by the system's first design principle.


%%%%%%%%%%%%%%%%%%%%%%%%%%%%%%%%%%%%%%%%%%%%%%%%%%%%
\subsection{Query-Dependent Authority Scoring}
\label{sec:authority-scoring}

\subsubsection{Positioning and Design Rationale}

Authority is a fundamentally relational concept: a speaker is authoritative with respect to
a topic, not in the abstract. This observation has direct algorithmic consequences. Static
authority measures---such as citation counts, aggregate legislative productivity, or tenure
length---can at best capture a form of institutional prestige that is orthogonal to the
question of domain expertise. A senator with twenty years of legislative experience on fiscal
policy is not, by virtue of that record alone, an authority on public health; yet a static
prestige measure would score them identically across both domains.

The system therefore models authority as a function of two arguments: the speaker's profile
$s$ and the query embedding $\mathbf{q}$. The composite authority score
$A(s, q) \in [0, 1]$ is the weighted sum of six components, each independently normalised to
the unit interval:

\begin{equation}
    A(s, q) = \sum_{k=1}^{6} w_k \cdot f_k(s, q), \quad \sum_{k=1}^{6} w_k = 1,
    \label{eq:authority}
\end{equation}

where each $f_k(s, q) \in [0, 1]$ captures a distinct dimension of the relationship between
the speaker's profile and the query's semantic domain. The weights $\{w_k\}$ are fixed
hyperparameters that encode the relative diagnostic value of each dimension.

The multi-factor design reflects the observation that no single dimension of expertise is
individually sufficient. Legislative activity is a strong signal of domain engagement but can
be inflated by co-signatorship on peripheral acts. Committee membership provides a structural
indicator of institutional assignment but does not capture self-directed specialisation outside
formal roles. Professional and educational background are the most stable signals but are also
the most distal from current parliamentary activity. The synthesis of six complementary signals
produces a score that is more robust to manipulation and more stable across queries than any
single-factor measure.

\subsubsection{Component Definitions}

\paragraph{Professional alignment.}
The profession component assesses the semantic proximity between the speaker's declared
professional background and the query domain. Each speaker $s$ is associated with a
pre-computed embedding $\mathbf{p}_s \in \mathbb{R}^d$ of their professional description.
The component is defined as:

\begin{equation}
    f_\text{prof}(s, q) = \max\!\left(0,\; \frac{\mathbf{p}_s \cdot \mathbf{q}}{\|\mathbf{p}_s\| \cdot \|\mathbf{q}\|}\right),
\end{equation}

where the ReLU clipping ensures that negative cosine similarities---which arise when the
professional background is semantically opposite to the query domain---are mapped to zero
rather than contributing negatively to the composite score. This choice reflects a conservative
principle: anti-expertise is not modeled as negative authority, only as the absence of a
positive signal.

\paragraph{Educational alignment.}
The education component is structurally identical to the profession component but operates on
a separately pre-stored embedding $\mathbf{u}_s \in \mathbb{R}^d$ of the speaker's educational
history:

\begin{equation}
    f_\text{educ}(s, q) = \max\!\left(0,\; \frac{\mathbf{u}_s \cdot \mathbf{q}}{\|\mathbf{u}_s\| \cdot \|\mathbf{q}\|}\right).
\end{equation}

Professional and educational backgrounds are treated as separate components rather than
concatenated into a single profile embedding because they carry distinct epistemological status:
educational formation reflects foundational disciplinary competence acquired prior to
parliamentary activity, whereas professional background reflects applied expertise developed
through practice. A physician who later became a businessman and then entered parliament
presents a mixed signal that concatenation would obscure.

\paragraph{Committee membership.}
The committee component evaluates the topical relevance of the speaker's formally assigned
parliamentary roles. Each committee $m$ in the system's thematic ontology is associated with
a set of domain-specific keywords whose joint embedding $\mathbf{k}_m \in \mathbb{R}^d$ serves
as a semantic descriptor of the committee's area of competence. The component selects the
maximally relevant active committee membership:

\begin{equation}
    f_\text{comm}(s, q) = \max_{m \in \mathcal{M}_s(t)}\; \max\!\left(0,\; \frac{\mathbf{k}_m \cdot \mathbf{q}}{\|\mathbf{k}_m\| \cdot \|\mathbf{q}\|}\right),
\end{equation}

where $\mathcal{M}_s(t)$ denotes the set of committees to which speaker $s$ holds a valid
membership at the reference time $t$ of the query. The max-over-memberships aggregation reflects
the principle that authority on a given topic is determined by the speaker's most relevant
institutional assignment, not by the average relevance of all their roles.

\paragraph{Parliamentary acts.}
The acts component aggregates the speaker's legislative activity by weighting each sponsored
act according to its topical relevance to the query and its temporal proximity to the reference
date:

\begin{equation}
    f_\text{acts}(s, q) = \sigma\!\left(\sum_{a \in \mathcal{A}_s} \rho_a \cdot \delta_a \cdot e^{-\lambda_\text{acts}(t - t_a)} \cdot \max\!\left(0,\; \frac{\mathbf{a} \cdot \mathbf{q}}{\|\mathbf{a}\| \cdot \|\mathbf{q}\|}\right)\right),
\end{equation}

where $\mathcal{A}_s$ is the set of acts co-signed by speaker $s$, $\rho_a \in (0, 1]$ is a
role coefficient that distinguishes primary sponsorship from co-signatorship, $\delta_a$ is a
relevance gate that sets the contribution to zero when the act's topical similarity to the
query falls below a minimum threshold, $\lambda_\text{acts}$ is the decay rate parameter for
legislative activity, and $\sigma(\cdot)$ is the logistic function, applied to map the unbounded
sum onto the unit interval. The threshold gate $\delta_a$ is essential to prevent a prolific
legislator whose acts span many domains from accumulating a high score through sheer volume
of off-topic activity.

\paragraph{Previous speeches.}
The interventions component applies the same aggregation logic to the speaker's speech history:

\begin{equation}
    f_\text{spch}(s, q) = \sigma\!\left(\sum_{v \in \mathcal{V}_s} \delta_v \cdot e^{-\lambda_\text{spch}(t - t_v)} \cdot \max\!\left(0,\; \frac{\mathbf{v} \cdot \mathbf{q}}{\|\mathbf{v}\| \cdot \|\mathbf{q}\|}\right)\right),
\end{equation}

where $\mathcal{V}_s$ is the set of speeches attributed to speaker $s$, $\mathbf{v}$ is the
pre-stored embedding of each speech, and $\lambda_\text{spch}$ is a decay rate that is set
strictly greater than $\lambda_\text{acts}$. The higher decay rate for speeches reflects the
greater temporal sensitivity of verbal interventions: parliamentary debates are responses to
current legislative and political events, and a speech that was highly pertinent two years ago
may be irrelevant or even misleading when applied to the current query context.

\paragraph{Institutional role.}
The role component assigns a baseline authority weight to speakers who occupy formal
institutional positions:

\begin{equation}
    f_\text{role}(s, q) = \alpha_r + \beta_r \cdot \max\!\left(0,\; \frac{\mathbf{r}_s \cdot \mathbf{q}}{\|\mathbf{r}_s\| \cdot \|\mathbf{q}\|}\right),
\end{equation}

where $\alpha_r \in [0, 1]$ is a fixed base weight determined by the institutional tier of the
role (committee president, vice-president, secretary, government member), $\beta_r$ is a scaling
coefficient for the semantic alignment bonus, and $\mathbf{r}_s$ is an embedding of the
speaker's role title. The base weight ensures that institutional authority is not wholly
contingent on the semantic alignment of a role title with the query---a minister of health
commands institutional authority on healthcare matters even if the alignment between the title
embedding and a highly technical query is imperfect.

\subsubsection{Temporal Decay}

The exponential decay factors in the acts and interventions components model the depreciation
of authority evidence over time. The general form is:

\begin{equation}
    \delta(t, t_0) = e^{-\lambda (t - t_0)},
\end{equation}

where $t$ is the reference date of the query, $t_0$ is the date of the parliamentary act or
speech, and $\lambda > 0$ is the decay rate parameter. The half-life of the decay function,
the time after which a piece of evidence retains half its original weight, is:

\begin{equation}
    h = \frac{\ln 2}{\lambda}.
\end{equation}

The choice of half-life is governed by the epistemological interpretation of what constitutes
relevant prior activity in each domain. Legislative commitments---acts and their signatorship
records---have a longer effective half-life because they reflect strategic policy positions
that tend to persist across multiple sessions and because the formal record of sponsorship
remains part of the deputy's political identity long after the act is voted upon. Verbal
interventions, by contrast, are responses to the current parliamentary agenda: a deputy may
speak urgently about a topic in one session that has been superseded by subsequent legislative
developments in the next.

Beyond the pragmatic differentiation between the two signal types, the temporal decay mechanism
serves a normative function: it prevents the authority score from being dominated by deputies
whose parliamentary prominence peaks during a period that predates the query by several
legislative cycles. In a parliament whose composition evolves over time---with deputies entering
and leaving, parties forming and dissolving, policy coalitions shifting---an authority measure
that does not discount historical activity is epistemically misleading.

\subsubsection{Coalition Logic}

Parliamentary group affiliation in the Italian legislature is not a static property. Deputies
may switch their group membership during a legislature, and the composition of coalition
alliances changes between legislative terms. This temporal dynamics of affiliation has direct
consequences for authority scoring: a deputy's authority score should be interpreted within the
context of their current political alignment, not their historical one.

The system resolves group affiliation at query time by querying the temporal validity attributes
of the membership relationship between a deputy and their parliamentary group. Formally, the
active group $g_s(t)$ of speaker $s$ at reference time $t$ is the unique group $g$ such that
the membership relation $(s, g)$ is valid at $t$---that is, $t_\text{start}(s, g) \leq t \leq
t_\text{end}(s, g)$. This resolution ensures that the authority score is associated with the
correct political identity of the speaker at the time the query is processed.

The coalition logic has a subtler implication for the aggregation of authority across groups.
When the system assigns the top-authority speaker for each parliamentary group to the expert
panel displayed alongside the answer, it must use the query-time group affiliation of each
speaker. A deputy who was affiliated with a centre-right group during a previous legislature
but has since joined a centrist formation should be credited to the current formation, not the
historical one. Failure to enforce this temporal validity condition would introduce a systematic
historical bias: the authority panel would reflect the political landscape of past legislatures
rather than the current session, misrepresenting the ideological positions of parliamentary
groups as they exist at the time of the query.

\subsubsection{Normalisation}

The six component scores are each independently normalised to $[0, 1]$ before being combined
by the weighted sum in Equation~\ref{eq:authority}. The normalisation is percentile-based
rather than range-based. For each component $f_k$ and each query, the raw scores of all speakers
in the population are ranked, and each speaker's normalised score is set to their percentile rank:

\begin{equation}
    \tilde{f}_k(s) = \frac{|\{s' \in \mathcal{S} : f_k(s', q) \leq f_k(s, q)\}|}{|\mathcal{S}|},
\end{equation}

where $\mathcal{S}$ is the full population of speakers in the corpus.

The choice of percentile normalisation over linear min-max scaling is motivated by the
distributional properties of the individual components. Several of them---most notably the
acts and interventions components---exhibit heavy-tailed distributions across the speaker
population: a small number of highly active deputies account for a disproportionate share of
the total legislative and discursive activity. Linear normalisation under such distributions
would compress the scores of the vast majority of deputies into a narrow range near zero,
rendering the composite score effectively determined by the outliers. Percentile normalisation
spreads the scores uniformly across the unit interval by construction, ensuring that the
composite score reflects relative standing within the population rather than raw magnitude.

A secondary benefit of percentile normalisation is query-independent comparability: because
the normalisation is applied over the full population for each query, scores from different
queries are expressed in the same distributional frame and can be compared meaningfully in
evaluation contexts. This property is essential for the evaluation pipeline, which must compare
the authority profiles of system-cited experts across a diverse set of evaluation queries.


%%%%%%%%%%%%%%%%%%%%%%%%%%%%%%%%%%%%%%%%%%%%%%%%%%%%
\subsection{Four-Stage Generation Pipeline}
\label{sec:generation-pipeline}

\subsubsection{Problem Definition}

Multi-view generation under parliamentary constraints can be formalised as the problem of
producing, from a query $q$ and a set of evidence fragments $\mathcal{R}$, a structured
response $\mathcal{O}$ that satisfies three simultaneous constraints.

The \textit{party coverage constraint} requires that, for each parliamentary group $g_j \in G$
represented in $\mathcal{R}$, the output $\mathcal{O}$ contains at least one section $s_{g_j}$
that articulates the position of group $g_j$ based exclusively on evidence attributed to its
members: $\forall g_j \in G,\; \exists\, s_{g_j} \subset \mathcal{O}$ such that every
evidentiary claim in $s_{g_j}$ is supported by some $c_i \in \mathcal{R}$ with
$\gamma(c_i) = g_j$.

The \textit{citation grounding constraint} requires that every claim in $\mathcal{O}$ that is
attributed to a specific speaker be supported by a verbatim quotation extracted from the
corresponding source evidence: $\forall$ attributed claim $x \in \mathcal{O}$, $\exists\, c_i
\in \mathcal{R}$ such that the quoted string $x_\text{quote}$ is a contiguous substring of the
raw text of $c_i$.

The \textit{coherence constraint} requires that the sections of $\mathcal{O}$, despite being
generated independently per group, be integrated into a discourse that is logically coherent
and free of internal contradiction at the meta-level---that is, the integrating discourse does
not misrepresent the positions of any group or introduce claims not supported by any evidence
in $\mathcal{R}$.

These three constraints are mutually incompatible for single-pass generation. A single
large-language-model invocation that receives the full evidence set and is asked to produce a
multi-view answer satisfying all three simultaneously faces an attention allocation problem: the
model cannot reliably maintain separate evidentiary chains for each group while simultaneously
enforcing citation placeholders and producing a coherent integrative discourse. Empirically,
single-pass generation tends to produce one of two failure modes: either the model merges the
positions of different groups into an unattributed synthesis that violates the coverage
constraint, or it produces correctly attributed sections that are individually faithful but
mutually incoherent when read as a unified document. The four-stage decomposition resolves
this tension by assigning each constraint to the stage best suited to enforce it.

\subsubsection{Stage 1: Claim Analyst}

The first stage receives the query $q$ and the full evidence set $\mathcal{R}$ and produces a
decomposition of the query into a set of \textit{atomic claims} $\mathcal{K} = \{k_1, \ldots,
k_L\}$, where each $k_l$ is an elementary propositional question that a complete answer to $q$
would need to address.

The purpose of this decomposition is to establish a shared epistemic scaffold that governs all
subsequent generation stages. Without a common claim structure, the party-specific sections
produced in Stage 2 would pursue independent narrative threads, making the integration step in
Stage 3 structurally harder: the integrator would need to simultaneously discover the shared
argumentative agenda and weave the sections together, which are tasks that are individually
demanding and jointly prone to failure. By externalising the claim structure as an explicit
intermediate representation, the pipeline separates the \textit{question of what to address}
from the \textit{question of how each party addresses it}.

The decomposition also serves a coverage verification function. A claim that is addressed by
some parties but absent from the section of another signals either a genuine gap in that party's
evidential record or a failure of the sectional writer to engage with the shared agenda. The
claim structure makes this distinction tractable: it becomes possible to audit, for each group
section, which claims are addressed and which are missing, and to intervene algorithmically if
the omission is due to generation failure rather than genuine absence of evidence.

\subsubsection{Stage 2: Sectional Writer}

The second stage generates, for each parliamentary group $g_j \in G$, an independent
narrative section $s_{g_j}$ that addresses the claim set $\mathcal{K}$ using only the evidence
fragments in $\mathcal{R}_{g_j} = \{c_i \in \mathcal{R} : \gamma(c_i) = g_j\}$.

The group-scoped evidence restriction is the primary mechanism for enforcing the party coverage
constraint and for preventing cross-group confabulation. If a single invocation were provided
with the full evidence set, the generator would have no structural obligation to distinguish
between evidence from different groups; it could, and empirically does, produce attributions
that blend positions across groups or assign evidence to the wrong speaker. By presenting each
invocation with only the evidence from a single group, the attribution problem is made
unambiguous: every claim the generator produces must be supported by the evidence before it,
and every piece of that evidence belongs to the same political formation.

Within each section, citations are not inserted as verbatim quotations but as structured
\textit{citation placeholders}. Each placeholder encodes the identifier of the source chunk
whose content the generator intends to cite, together with an inline approximation of the
intended quotation extracted from the evidence text. This representation serves two purposes.
It defers the verbatim extraction to the deterministic Stage 4, which can perform it without
LLM involvement, and it provides a traceable record of \textit{which} evidence the generator
considered relevant for \textit{which} claim. The collection of all placeholders across all
sections constitutes the \textit{citation registry}, which is the authoritative record of
citation expectations against which later stages are audited.

\subsubsection{Stage 3: Narrative Integrator}

The third stage receives the collection of party-specific sections $\{s_{g_1}, \ldots,
s_{g_M}\}$ and synthesises them into a unified discourse $\mathcal{O}_\text{int}$ that
reads as a coherent document while preserving the attribution and evidential basis of each
section.

The coherence constraint is inherently in tension with the structural autonomy of the sectional
writing stage. Sections generated independently may use different terminological conventions,
address the same claim at different levels of abstraction, or organise their arguments in
incompatible narrative orders. The integrator must resolve these surface inconsistencies without
altering the substantive content of any section and without suppressing the citation placeholders
that bind each claim to its source.

The citation preservation requirement is enforced through a verification mechanism that audits
the placeholder count in $\mathcal{O}_\text{int}$ against the citation registry established in
Stage 2. If the count in the integrated output is lower than the registry count, citations have
been dropped during synthesis. In such cases, a repair procedure reinjests the content of the
affected sections as fallback paragraphs appended to the integrated discourse, ensuring that no
evidence binding established in Stage 2 is silently invalidated by the integration process.

The integrator operates under an explicit instruction to avoid producing claims that are not
grounded in any section---that is, it is not permitted to use the synthesis step as an
opportunity to introduce new propositional content unsupported by the sectional evidence. This
constraint is not enforceable deterministically, but it is systematically reinforced by the
citation verification mechanism: any new claim inserted during integration would either carry
no citation placeholder or would be paired with a placeholder that does not appear in the
registry, and both cases are detectable.

\subsubsection{Stage 4: Citation Surgeon}

The fourth stage is not a generative operation but a deterministic extraction procedure. Its
function is to resolve each citation placeholder in $\mathcal{O}_\text{int}$ into a verified
verbatim quotation drawn from the raw parliamentary record, replacing the structured placeholder
with the extracted text.

For each placeholder, the surgeon retrieves from the data layer the byte-level character
offsets $(b_\text{start}, b_\text{end})$ associated with the source chunk identified in the
placeholder. It then extracts the contiguous character span
$\text{raw}[b_\text{start} : b_\text{end}]$ from the unmodified original speech text, using
the alignment map to translate between the cleaned representation used by the embedding and
retrieval stages and the authoritative original source used for citation.

The inline approximation embedded in the placeholder by the sectional writer is treated as a
candidate quotation that must pass a faithfulness verification before acceptance. The surgeon
checks whether the inline string is a literal substring of the extracted raw span. If the check
passes, the inline string is used as the citation text, which may be a meaningful sub-span of
the full chunk. If the check fails---which can occur when the generator has paraphrased or
truncated the source passage during section writing---the surgeon falls back to a
semantically-guided extraction that scores individual sentences within the chunk span by
their cosine similarity to the query embedding and selects the highest-scoring sentence as the
citation text. In both cases, the resulting quotation is a verbatim substring of the original
source document.

This stage embodies the \textit{citation faithfulness invariant} as a computational guarantee
rather than a normative aspiration. The guarantee is possible precisely because citation text is
never produced by a generative language model: the model may select which evidence to cite, but
the text of the citation is always determined by a deterministic lookup into the archived source.
The structural consequence is that no failure mode of the generative stages can produce a
citation that is a paraphrase, a confabulation, or a semantically approximate substitution for
the original; the worst failure is a dropped citation, which the registry verification mechanism
in Stage 3 is designed to detect and repair.


%%%%%%%%%%%%%%%%%%%%%%%%%%%%%%%%%%%%%%%%%%%%%%%%%%%%
\subsection{Citation Integrity System}
\label{sec:citation-integrity}

The citation integrity system is a cross-cutting component that operates in parallel with the
generation pipeline rather than after it. Its function is to maintain the binding between
textual claims and their evidential sources throughout the multi-stage process, preventing
the progressive loss of attribution information that characterises monolithic generation
approaches and providing the structural substrate for the faithfulness guarantee described in
the previous section.

\subsubsection{Citation Registry}

The citation registry is a persistent data structure that accumulates citation expectations
as the generation pipeline progresses. It is initialised at the beginning of Stage 2 and
populated incrementally as each party-specific section is produced. For each citation
placeholder $p_l$ emitted by the sectional writer, the registry records the triple
$(p_l,\, \text{id}(c_{i(l)}),\, \mathbf{q})$, where $p_l$ is the placeholder identifier,
$\text{id}(c_{i(l)})$ is the unique identifier of the source chunk referenced in the
placeholder, and $\mathbf{q}$ is the query embedding, which will be used in the fallback
extraction step if the inline approximation fails verification.

The registry is \textit{append-only} during Stage 2: entries are added but never deleted.
This design ensures that the complete history of citation expectations established during
sectional writing is preserved regardless of what subsequent stages do to the text. The
verification step in Stage 3 consults this registry to detect placeholder deletion, and the
extraction step in Stage 4 uses it to resolve each placeholder to its source coordinates.

The registry serves an additional auditing function during evaluation: by comparing the
registered citation expectations with the citations that survive to the final output, it
is possible to compute a citation retention rate that quantifies the faithfulness of the
generation pipeline as a whole. This rate is one of the automated metrics used in the
evaluation framework described in the following chapter.

\subsubsection{Citation-First Architecture}

The design philosophy of the citation integrity system inverts the conventional relationship
between generation and attribution. In standard retrieval-augmented generation, the language
model produces text first and citations are appended or verified afterwards---a post-hoc
attribution model in which the generator is free to produce any claim and citations are
sought to support it. This model is structurally susceptible to confabulation: a confident
generator may produce plausible-sounding claims for which no supporting evidence exists, and
the post-hoc attribution step may then match those claims to superficially related but
actually unsupportive passages.

The citation-first architecture reverses this logic. The evidence fragments are fixed before
any generation begins; the generator is constrained to write \textit{around} the evidence,
inserting placeholders at the point of citation rather than generating text and then
annotating it. The consequence is that every claim in the output is, by construction, linked
to an evidence fragment that existed in the retrieved set before the generation stage was
invoked. The generator cannot introduce evidentiary claims that lack a corresponding source
entry in the citation registry, because the registry is populated from the evidence presented
to the generator and any placeholder that does not correspond to an existing evidence fragment
will be rejected during the extraction stage.

This inversion is not merely a procedural constraint; it has a semantic consequence for the
generation task itself. When the sectional writer is presented with a set of evidence
fragments and instructed to produce a narrative that cites them, it is performing a task closer
to \textit{guided selection and arrangement} than to \textit{open-ended synthesis}. The
generator must decide which evidence is most relevant to each atomic claim, in what order to
present it, and how to connect citations into a coherent argument---but it cannot decide what
the evidence says. This narrowing of the generation task reduces the surface area for
hallucination and keeps the epistemic authority of the output grounded in the original
parliamentary record.

\subsubsection{Coherence Validation}

The citation integrity system includes a semantic validation layer that operates on the final
output to verify that the relationship between each cited passage and its surrounding claim is
coherent---that is, that the claim is genuinely supported by the cited text and not merely
proximate to it.

For each citation in the final output, the validation layer computes the cosine similarity
between the embedding of the surrounding claim context and the embedding of the cited passage:

\begin{equation}
    \text{coh}(x, c_i) = \frac{\phi(x) \cdot \mathbf{e}_i}{\|\phi(x)\| \cdot \|\mathbf{e}_i\|},
\end{equation}

where $\phi(x)$ is the embedding of the claim context $x$ and $\mathbf{e}_i$ is the embedding
of the cited chunk. A coherence score below a minimum threshold $\tau_\text{coh}$ indicates
that the citation is semantically mismatched with its enclosing claim---a situation that can
arise when the integrator reorganises sections in a way that displaces a citation from the
claim it was originally intended to support.

Mismatched citations detected by the validation layer are flagged in the output metadata rather
than silently corrected, for two reasons. First, the correct resolution of a mismatch---whether
to reassign the citation, to delete the claim, or to reorder the surrounding text---requires
a judgment about argumentative intent that cannot be made deterministically. Second, making
the mismatch visible in the metadata preserves the transparency of the system's output: a
downstream user who examines the citation provenance can identify cases where the system
flagged a coherence concern and evaluate the cited passage independently. Opacity in citation
quality is a more serious failure mode for a system designed to support journalistic and civic
accountability than a conservative flagging of borderline cases.

The combination of the citation registry, the citation-first generation architecture, and the
coherence validation layer constitutes a three-layered defence of the citation faithfulness
invariant. The registry prevents citation loss by maintaining a complete record of expectations.
The citation-first architecture prevents confabulation by structuring generation around
pre-existing evidence. The coherence validation prevents semantic displacement by verifying that
cited evidence is genuinely relevant to the claims it accompanies. No single layer is sufficient
alone; together, they provide structural guarantees that are substantially stronger than any
post-hoc verification approach applied to a monolithic generation process.
